\section{Methodology}
\label{sec:method}
The following equations describe the active branches of the frozen implementation. In particular, the final evolution changes the calibration distance, not the covariance acquisition formula. Unused alternative losses remaining in the source are not components of the reported method.

\subsection{Compatible mechanisms and calibrated weights}
Let $\bar X_{metc}=S^{-1}\sum_sX_{msetc}$ and let $H_{etc}$ be the observed historical trajectory mean. For the observed history records $\mathcal O$, define separate process and trajectory discrepancies
\begin{align}
 B_m&=|\mathcal O|^{-1}\sum_{r\in\mathcal O}(p_{mr}-y_r)^2,\nonumber\\
 R_m&=\left[(ETC)^{-1}\sum_{etc}(\bar X_{metc}-H_{etc})^2\right]^{1/2}.
 \label{eq:compaterrors}
\end{align}
When $\mathcal O$ is empty, $B_m=0$. The compatible set is
\begin{equation}
 \mathcal C=\{m:B_m\leq\min_j B_j+0.045\}\cap\{m:R_m\leq0.10\}.
 \label{eq:set}
\end{equation}
An empty set raises an error instead of silently widening tolerances. These empirical gates provide neither confidence-set coverage nor a guarantee that the true mechanism is on the grid. In particular, process Brier discrepancy is a screening statistic, not a Bernoulli log likelihood.

The final calibration rule normalizes trajectory residuals by estimated simulation noise and emphasizes the worst historical event:
\begin{align}
 v_{etc}&=\frac1{MS}\sum_{m,s}(X_{msetc}-\bar X_{metc})^2,
 &\bar v&=\frac1{ETC}\sum_{etc}v_{etc},\nonumber\\
 r_{etc}&=\frac{v_{etc}+\epsilon}{\bar v+\epsilon},
 &D_m&=\max_e\frac1{TC}\sum_{tc}\frac{(\bar X_{metc}-H_{etc})^2}{r_{etc}},\label{eq:distance}\\
 w_m^{(0)}&=\frac{\ind\{m\in\mathcal C\}\exp[-D_m/(2\sigma^2)]}
 {\sum_{j\in\mathcal C}\exp[-D_j/(2\sigma^2)]},
 &\epsilon&=10^{-5},\quad\sigma=7\times10^{-4}.\label{eq:weights}
\end{align}
Variance is computed with divisor $S$ and averaged over all $M$ candidates before screening. It is seed variance, not variance of the Monte Carlo mean. This supplies time- and channel-dependent precision weights through $r_{etc}$; there is no separately learned temporal attention term or higher-order acquisition correction. The resulting distribution is a Gibbs calibration mixture. Calling it an exact posterior for the historical trajectory-generating process would require an observation likelihood that the implementation does not specify.

\subsection{Missing viewing records and observation likelihoods}
In the standardized audit cohort, unobserved viewing is integrated out. If $V_r$ indicates viewing and $I_r$ an intention, then
\begin{equation}
 P(I_r=1\mid\theta_m)=P(V_r=1\mid\theta_m)P(I_r=1\mid V_r=1,\theta_m)
 =a_m\rho_{e(r)}b_m\iota_r,
 \label{eq:view}
\end{equation}
where $\rho_e$ is notification reach and $\iota_r$ is recorded interest. Thus a missing view is not a negative view. Waiting-dropout records condition on a known waiting-risk state. These formulas are not substitutes for filtering an unobserved interaction history.

For a revealed record with observation indicator $o_r$, the supplied MCAR benchmark uses $\ell_{mr}=p_{mr}^{y_r}(1-p_{mr})^{1-y_r}$ when $o_r=1$ and $\ell_{mr}=1$ otherwise. The optional known observation model $\gamma_{mry}=P(o_r=1\mid Y_r=y,\theta_m)$ instead gives
\begin{equation}
 \ell_{mr}=\begin{cases}
 p_{mr}^{y_r}(1-p_{mr})^{1-y_r}\gamma_{mr y_r},&o_r=1,\\
 (1-p_{mr})(1-\gamma_{mr0})+p_{mr}(1-\gamma_{mr1}),&o_r=0.
 \end{cases}
 \label{eq:missing}
\end{equation}
This distinction follows the need to model nonignorable observation processes \citep{rubin1976}. The optional branch is unit-tested, not empirically validated under real missingness. Right-censoring is different: a pending time-to-event observation would require a survival likelihood such as $P(T_r>t_{\mathrm{cut}}\mid\theta_m)$ and a specified event-time model. That extension is not implemented; our evidence pool consists of completed binary opportunities.

\subsection{Cached probes and the covariance acquisition identity}
Let $z_m\in\R^D$ concatenate the candidate's mean trajectories on the fixed combination probes, with $D=E_{\mathrm p}TC$. Given current weights $w$, define
\begin{equation}
 \mu=\sum_mw_mz_m,\qquad q_r=\sum_mw_mp_{mr},\qquad
 c_r=\sum_mw_m(p_{mr}-q_r)(z_m-\mu).
 \label{eq:cov}
\end{equation}
The acquisition value is
\begin{equation}
 A_r(w)=\frac{\|c_r\|_2^2}{Dk_rq_r(1-q_r)},\qquad
 C=P^\top\operatorname{diag}(w)(Z-\mathbf1\mu^\top).
 \label{eq:score}
\end{equation}
Rows of $C$ contain all $c_r^\top$. Computing the matrix product shares work across records. The denominator is clamped below by $10^{-12}$ in code; deterministic observations have zero exact value.

\begin{proposition}[Binary predictive-variance reduction]
For fixed cached conditional means $Z=z_\theta$, the model-implied binary observation law, and $0<q_r<1$, the expected reduction in mean coordinate variance after observing $Y_r$ equals $k_r A_r(w)$.
\end{proposition}
\begin{proof}
Let $\mu_y=\E[Z\mid Y_r=y]$. Bayes' rule yields $\mu_1-\mu=c_r/q_r$ and $\mu_0-\mu=-c_r/(1-q_r)$. Applying total covariance gives
\begin{align}
 \frac1D\tr\{\Cov(Z)-\E_{Y_r}\Cov(Z\mid Y_r)\}
 &=\frac{q_r\|\mu_1-\mu\|^2+(1-q_r)\|\mu_0-\mu\|^2}{D}\nonumber\\
 &=\frac{\|c_r\|^2}{Dq_r(1-q_r)}.
\end{align}
\end{proof}
This is a specialization of predictive variance reduction, not a new general optimal-design theorem. Its role is to justify an efficient implementation. It requires that the observation change the target only through mechanism weights. If obtaining a record intervenes on an agent or reveals a target-correlated latent state not represented by $\theta$, the conditional trajectories must also change; the cached formula alone is then insufficient. Nor does internal expected variance reduction imply a reduction in realized RMSE, aleatoric variance, or model discrepancy.

\subsection{Small-batch conditioning and predictive uncertainty}
To avoid evaluating every member of a batch under the same unupdated weights, maintain hypothetical branches $\mathcal F_j=\{(\pi_f,w_f)\}$ after $j$ selections. The next record maximizes
\begin{equation}
 r_{j+1}=\argmax_{r\ \mathrm{eligible}}\sum_{f\in\mathcal F_j}\pi_f A_r(w_f).
 \label{eq:batch}
\end{equation}
Each branch splits under $y\in\{0,1\}$ using $L_{mry}=p_{mr}^y(1-p_{mr})^{1-y}$, branch probability $q_{fy}=\sum_mw_{fm}L_{mry}$, mass $\pi_fq_{fy}$, and weights $w_{fm}L_{mry}/q_{fy}$. Selection reads no realized labels. This exactly averages hypothetical outcomes for a greedy batch under conditional independence, while branches of probability at most $10^{-12}$ are omitted numerically. It is neither globally optimal batch search nor a generally submodular objective; Section~\ref{sec:boundaries} gives a counterexample. Batch size is two, with an implementation limit of four.

After acquisition, the actual evidence updates weights according to
\begin{equation}
 w'_m\propto w_m\exp\left\{\tau^{-1}\sum_{r\in Q}\log\ell_{mr}\right\},\qquad\tau=1.
 \label{eq:update}
\end{equation}
Logarithm inputs are clipped to $[10^{-7},1-10^{-7}]$; zero prior mass remains zero. The target distribution is a mixture of complete cached paths,
\begin{equation}
 \widehat P(X^*\mid Y_Q)=\sum_{m,s}\frac{w'_m}{S}\delta_{X^*_{ms}},
 \qquad\widehat\mu=\sum_mw'_m\bar X^*_m.
 \label{eq:mixture}
\end{equation}
Its coordinatewise uncertainty decomposes as
\begin{align}
 V_{\mathrm{epi}}&=\sum_mw'_m(\bar X^*_m-\widehat\mu)^2,\nonumber\\
 V_{\mathrm{alea}}&=\sum_mw'_m\frac1S\sum_s(X^*_{ms}-\bar X^*_m)^2.
 \label{eq:uncertainty}
\end{align}
Weighted fifth and ninety-fifth percentiles give pointwise 90\% predictive intervals. They are not simultaneous path bands or confidence intervals for the real population. Complete atoms preserve temporal dependence even when marginal plots do not display it.

\begin{algorithm}[tbp]
\caption{Compatible-mechanism inference with cached process evidence}
\label{alg:main}
\small
\begin{algorithmic}[1]
\Require Historical data; caches $X,P,Z,X^*$; eligible pool $U$; costs $k$; budget $B$
\State Compute $B_m,R_m$ and $\mathcal C$ by Eqs.~\eqref{eq:compaterrors}--\eqref{eq:set}
\If{$\mathcal C=\emptyset$}
 \State \textbf{raise} model-inadequacy error
\EndIf
\State Compute $v,r,D$; initialize $w$ by Eqs.~\eqref{eq:distance}--\eqref{eq:weights}
\While{an affordable record remains in $U$}
 \State $\mathcal F\gets\{(1,w)\}$, $Q\gets\emptyset$, $B_{\mathrm{reserve}}\gets B$
 \For{$j=1,2$}
  \State $U'\gets\{r\in U\setminus Q:k_r\leq B_{\mathrm{reserve}}\}$
  \If{$U'=\emptyset$} \State \textbf{break} \EndIf
  \State Score all $r\in U'$ using Eqs.~\eqref{eq:score} and \eqref{eq:batch}
  \State Select maximum-score $r$; add to $Q$; reserve $k_r$
  \State Expand each branch for both binary outcomes of $r$
 \EndFor
 \State Retrieve $Y_Q$ and observation mask; compute Eq.~\eqref{eq:missing}
 \State Update $w$ by Eq.~\eqref{eq:update}; $U\gets U\setminus Q$, $B\gets B-\sum_{r\in Q}k_r$
\EndWhile
\State \Return whole-path mixture, mean, $V_{\mathrm{epi}}$, $V_{\mathrm{alea}}$, marginal quantiles
\end{algorithmic}
\end{algorithm}

\paragraph{Computational cost.}
Given $M$ mechanisms, $R$ records and $D$ probe coordinates, cached probabilities and conditional means require $O(MR+MD)$ memory. Full seed trajectories add $O(MSD)$ storage per bank. A covariance pass costs $O(MRD)$ arithmetic; explicit binary-outcome enumeration has the same order, but repeats centering and reductions in many small operations. The measured advantage is vectorization and shared computation, not a lower asymptotic complexity. Exact branch averaging for batch size $b$ costs $O((2^b-1)MRD)$ in the worst case, explaining the small-batch restriction. A label update costs $O(Mb)$, while target means and variances cost $O(MSD_*)$ and marginal sorting adds $O(D_*MS\log(MS))$. Offline simulation cost is paid before these operations. The mathematical interface is domain-independent, but both observation probabilities and the simulator bank remain domain-specific.
